\section{Related Work}

\subsection{Deep Learning for Cervical Cytology}
Automated cervical-cytology screening has been propelled by curated benchmarks, most prominently the SIPaKMeD database of Papanicolaou-stained conventional-smear cells annotated into five cytomorphological classes \citep{plissiti2018sipakmed}, and the earlier Herlev collection of single cells labelled under the WHO dysplasia schema \citep{jantzen2005herlev}. A large body of work reports near-ceiling within-dataset accuracy on these corpora using deep convolutional backbones, and recent liquid-based resources such as HiCervix extend the task to hierarchical Bethesda categories at substantially greater scale \citep{cai2024hicervix,nayar2015bethesda}. Convolutional architectures of the EfficientNet family \citep{tan2019efficientnet}, typically initialised from ImageNet pre-training \citep{deng2009imagenet}, have become a standard high-accuracy, parameter-efficient choice for this domain. However, the overwhelming majority of these studies evaluate on a held-out split of the same database on which they were trained, leaving open the question of whether the reported performance reflects transferable cytological competence or dataset-specific regularities.

\subsection{Knowledge Distillation for Efficient Medical Models}
The deployment of cytology classifiers on resource-constrained or point-of-care hardware motivates model compression, and knowledge distillation \citep{hinton2015distilling} has emerged as a leading strategy: a compact student is trained to match the temperature-softened logits of a larger teacher, frequently recovering---and occasionally exceeding---teacher accuracy through the regularising effect of soft targets. While distillation is now common across medical-imaging tasks, its interaction with out-of-distribution behaviour is rarely examined; a smaller, better-regularised student may generalise no better than its teacher if both have learned the same dataset-specific shortcuts.

\subsection{Generalization, Shortcut Learning, Calibration, and Explainability}
A recurring finding across medical imaging is that high in-distribution accuracy degrades sharply under cross-dataset evaluation, a phenomenon attributable in large part to shortcut learning, whereby networks exploit spurious correlations---staining, scanner, and preparation cues---rather than the intended morphology \citep{geirhos2020shortcut}. Two complementary tools are routinely invoked to interrogate such models. Confidence calibration, quantified by the Expected Calibration Error and reliability diagrams \citep{guo2017calibration}, assesses whether predicted probabilities are trustworthy, yet calibration established in-distribution is seldom verified to transfer across domains. Post-hoc attribution methods such as Grad-CAM \citep{selvaraju2017gradcam} localise the image regions driving a decision and are widely used to argue for clinical plausibility, although they reveal where a model attends rather than whether its reasoning is faithful.

\subsection{The Gap: Standardised Cross-Dataset Label Alignment}
Cross-dataset benchmarking in cervical cytology is further obstructed by incompatible taxonomies: the five SIPaKMeD morphology classes, the seven-grade Herlev dysplasia scheme, and the hierarchical Bethesda system of HiCervix and APACC do not share a common label space, and the requisite projection is frequently one-to-many or undefined \citep{nayar2015bethesda}. Prior comparisons rarely document the mapping or report the fraction of cases excluded because they have no valid correspondence (for instance, glandular lesions with no squamous home), which undermines reproducibility. Our work addresses this gap with a deterministic, fully documented binary Normal-versus-Abnormal alignment and an honesty metric quantifying the excluded fraction, enabling a leakage-aware audit of distilled-model generalisation across three independent Pap datasets.